\subsection{方程汇总}
\label{sec:chap1-equation-summary}

考虑黏度系数 $\lambda,\mu$ 为常数的 Newtonian 流体. 黏性应力张量为
\[
  T=2\mu D(v)+\lambda(\operatorname{div}v)\operatorname{Id},
\]
其中 $\lambda,\mu$ 满足条件\eqref{eq:chap1-viscosity-conditions}. 利用 Remark~\ref{rem:chap1-strain-rate-identities}, 流动由下列方程控制.

\begin{itemize}
\item 质量守恒:
\[
  \frac{\partial\rho}{\partial t}+\operatorname{div}(\rho v)=0.
\]
\item 线动量演化:
\[
\frac{\partial(\rho v)}{\partial t}+\operatorname{div}(\rho v\otimes v)-\mu\Delta v
+\nabla\bigl(p-(\lambda+\mu)\operatorname{div}v\bigr)-\rho f=0.
\]
动能方程可由前两个方程推出, 不必另加到流动演化模型中. 本书后文研究流体力学方程的数学方法均以此性质为基础, 通常称为能量方法.
\item 系统总能量的演化. 依所选热力学未知量不同, 该方程有多种形式上等价的写法, 每一种又有守恒和非守恒形式:
\begin{itemize}
\item 对单位质量总能量 $E=e+\tfrac12|v|^2$,
\[
\frac{\partial(\rho E)}{\partial t}+\operatorname{div}(\rho Ev)-\operatorname{div}(k\nabla T)
-\operatorname{div}(T\cdot v)+\operatorname{div}(pv)-\rho f\cdot v=0.
\]
\item 对内能,
\[
\frac{\partial(\rho e)}{\partial t}+\operatorname{div}(\rho ev)-\operatorname{div}(k\nabla T)
+p\operatorname{div}v-\lambda(\operatorname{div}v)^2-2\mu|D(v)|^2=0.
\]
\item 对比熵,
\[
\rho T\left(\frac{\partial s}{\partial t}+v\cdot\nabla s\right)-\operatorname{div}(k\nabla T)
-\lambda(\operatorname{div}v)^2-2\mu|D(v)|^2=0.
\]
\item 对温度, 并假定定容比热容 $c_v$ 为常数,
\[
\frac{\partial(\rho c_vT)}{\partial t}+\operatorname{div}(\rho c_vTv)-\operatorname{div}(k\nabla T)
+T\left(\frac{\partial p}{\partial T}\right)_\rho\operatorname{div}v
-\lambda(\operatorname{div}v)^2-2\mu|D(v)|^2=0.
\]
\end{itemize}
\end{itemize}

因此, 三个演化方程完全确定运动. 当然, 还必须加入一个把系统中各热力学变量联系起来的状态方程, 如理想气体定律 $p=k\rho T$, Van der Waals 定律 $(p+a\rho^2)(1-b\rho)=k\rho T$, 或其他状态方程, 以使系统封闭. 这组方程称为完整 Navier--Stokes--Fourier 系统. 还必须配以初始条件和边界条件, Section~\ref{subsec:chap1-initial-boundary-data} 将作简要讨论.

最后, 若与流动相关的所有场, 即密度, 压力, 速度和能量, 都与时间 $t$ 无关, 则称流动是稳态的. 但这绝不意味着流体静止, 后者要求 $v=0$.
